\documentclass{article}
\usepackage{amsmath, amssymb, amsthm}
\usepackage{hyperref}
\usepackage{geometry}
\geometry{margin=1in}

\title{Erd\H{o}s Problem \#14}
\date{Last updated: 2025-08-31}
\author{Source: erdosproblems.com}

\begin{document}
\maketitle

\noindent\textbf{Status:} OPEN \quad \textbf{No prize}

\noindent\textbf{Tags:} number theory, sidon sets, additive combinatorics

\medskip
\noindent\textbf{Problem Statement:}

Let $A\subseteq \mathbb{N}$. Let $B\subseteq \mathbb{N}$ be the set of integers which are representable in exactly one way as the sum of two elements from $A$.

Is it true that for all $\epsilon>0$ and large $N$\[\lvert \{1,\ldots,N\}\backslash B\rvert \gg_\epsilon N^{1/2-\epsilon}?\]Is it possible that\[\lvert \{1,\ldots,N\}\backslash B\rvert =o(N^{1/2})?\]

\bigskip
\noindent\textbf{Remarks:}

Apparently originally considered by Erd\H{o}s and Nathanson, although later Erd\H{o}s attributes this to Erd\H{o}s, S\'{a}rk\"{o}zy, and Szemer\'{e}di (but gives no reference), and claims a construction of an $A$ such that for all $\epsilon>0$ and all large $N$\[\lvert \{1,\ldots,N\}\backslash B\rvert \ll_\epsilon N^{1/2+\epsilon},\]and yet there for all $\epsilon>0$ there exist infinitely many $N$ where\[\lvert \{1,\ldots,N\}\backslash B\rvert \gg_\epsilon N^{1/3-\epsilon}.\]Erd\"{o}s and Freud investigated the finite analogue in \cite{ErFr91}, proving that there exists $A\subseteq \{1,\ldots,N\}$ such that the number of integers not representable in exactly one way as the sum of two elements from $A$ is $<2^{3/2}N^{1/2}$, and suggest the constant $2^{3/2}$ is perhaps best possible.

\bigskip
\noindent\textbf{References:}

[ErFr91] Erd\H{o}s, P. and Freud, R., On sums of a Sidon-sequence. J. Number Theory (1991), 196--205.

\bigskip
\noindent\small{Source: \url{https://www.erdosproblems.com/14}}
\end{document}
